\chapter{Mapa lógico de los dos temas}

\section{La idea central}

Sea $\A$ una álgebra de Banach compleja unital y $a\in\A$. La Parte 1 parte de una pregunta:
¿cuándo existe $(\lambda1-a)^{-1}$? La serie de Neumann responde cerca de un invertible; el
resolvente organiza la respuesta para todos los escalares; el espectro registra los valores donde
falla la inversión. La Sección 2 añade la identidad $C^*$ y obtiene consecuencias que no valen en
una álgebra de Banach arbitraria.

\section{Mapa de la Parte 1}

\begin{center}
\begin{tikzpicture}[node distance=9mm and 8mm]
  \node[theorybox,text width=3.25cm] (ban) {Álgebra de Banach unital\\completitud y producto};
  \node[theorybox,text width=3.25cm,right=of ban] (neu) {Serie de Neumann\\$\|x\|<1\Rightarrow(1-x)^{-1}$};
  \node[theorybox,text width=3.25cm,right=of neu] (open) {$\GL(\A)$ abierto\\inversión continua};
  \node[theorybox,text width=3.25cm,below=of neu] (bound) {Cota para $|\lambda|>\|a\|$\\espectro compacto};
  \node[theorybox,text width=3.25cm,right=of bound] (res) {Resolvente\\identidad y holomorfía};
  \node[theorybox,text width=3.25cm,below=of bound] (non) {Espectro no vacío\\Liouville + Hahn--Banach};
  \draw[arr] (ban)--(neu);
  \draw[arr] (neu)--(open);
  \draw[arr] (open)--(res);
  \draw[arr] (neu)--(bound);
  \draw[arr] (open.south west)--(bound.north east);
  \draw[arr] (res.south) |- (non.east);
  \draw[arr] (bound)--(non);
\end{tikzpicture}
\end{center}

La flecha que llega a la no vacuidad necesita dos ingredientes distintos: la holomorfía del
resolvente y su decaimiento en infinito, obtenido con la serie de Neumann. La compacticidad usa
la apertura de $\GL(\A)$ para el cierre y la misma serie para la acotación.

\begin{center}
\begin{tikzpicture}[node distance=5mm and 8mm]
  \node[theorybox,text width=3.25cm] (non) {Espectro no vacío};
  \node[theorybox,text width=3.25cm,right=of non] (gm) {Gelfand--Mazur\\álgebras de división};
  \node[theorybox,text width=3.25cm,right=of gm] (max) {Ideales maximales\\$\A/J\cong\mathbb C$ si $\A$ es conmutativa};
  \node[theorybox,text width=3.25cm,below=of max] (char) {Caracteres\\$\Omega(\A)$};
  \node[theorybox,text width=3.25cm,left=of char] (gel) {Transformada de Gelfand\\$a\mapsto\widehat a$};
  \node[theorybox,text width=3.25cm,left=of gel] (rad) {Radio espectral\\resolvente + potencias};
  \draw[arr] (non)--(gm);
  \draw[arr] (gm)--(max);
  \draw[arr] (max)--(char);
  \draw[arr] (char)--(gel);
  \draw[arr] (non)--(rad);
\end{tikzpicture}
\end{center}

El radio espectral ocupa una rama propia: su fórmula utiliza la holomorfía del resolvente,
la relación espectral para potencias y el principio de acotación uniforme; no se deduce de la
transformada de Gelfand. La identificación entre caracteres e ideales maximales requiere
conmutatividad, mientras que la teoría inicial del espectro no la requiere.

\section{Mapa de la Sección 2}

\begin{center}
\begin{tikzpicture}[node distance=5mm and 8mm]
  \node[theorybox,text width=3.25cm] (id) {Identidad $C^*$\\$\|a^*a\|=\|a\|^2$};
  \node[theorybox,text width=3.25cm,right=of id] (nr) {Para $a$ normal\\$\|a\|=r(a)$};
  \node[theorybox,text width=3.25cm,right=of nr] (uq) {Unicidad de la norma $C^*$\\por el radio de $a^*a$};
  \draw[arr] (id)--(nr);
  \draw[arr] (nr)--(uq);
\end{tikzpicture}
\end{center}

\begin{center}
\begin{tikzpicture}[node distance=5mm and 8mm]
  \node[theorybox,text width=3.25cm] (a) {$a$ normal en una $C^*$-álgebra unital};
  \node[theorybox,text width=3.25cm,right=of a] (comm) {$C^*(a,1)$ conmutativa\\por $aa^*=a^*a$};
  \node[theorybox,text width=3.25cm,right=of comm] (gn) {Gelfand--Naimark\\conmutativo};
  \node[theorybox,text width=3.25cm,below=of gn] (fc) {Cálculo funcional continuo\\$f\mapsto f(a)$};
  \node[theorybox,text width=3.25cm,left=of fc] (class) {Clasificación espectral\\y raíces positivas};
  \draw[arr] (a)--(comm);
  \draw[arr] (comm)--(gn);
  \draw[arr] (gn)--(fc);
  \draw[arr] (fc)--(class);
\end{tikzpicture}
\end{center}

La igualdad $\|a\|=r(a)$ requiere normalidad y la identidad $C^*$.
La conmutatividad de $C^*(a,1)$ se deduce directamente de
$aa^*=a^*a$. La transformada de Gelfand de la Parte 1 y la igualdad de
normas permiten la representación isométrica conmutativa. La permanencia
espectral identifica el espectro de $a$ en $C^*(a,1)$ con el espectro
ambiente. La raíz positiva y la descomposición polar usan después el cálculo
funcional, con las hipótesis de sus enunciados.

\section{Tres rutas de lectura}

\noindent\textbf{Invertibilidad y espectro.} Serie de Neumann $\Rightarrow$ estabilidad de los invertibles
$\Rightarrow$ resolvente holomorfo; junto con la estimación en infinito y Liouville, se obtiene
$\sigma_{\A}(a)\neq\varnothing$. En operadores $A\in\mathcal B(X)$, el mismo resolvente es
$(\lambda I_X-A)^{-1}$.

\noindent\textbf{Conmutatividad y funciones.} Para $\A$ conmutativa y unital, los ideales maximales
producen caracteres; éstos identifican $\sigma_{\A}(a)$ con el rango de $\widehat a$. Si $\A$
es además una $C^*$-álgebra, la representación conmutativa permite el cálculo funcional.

\noindent\textbf{Cocientes y operadores.} Para un espacio de Banach $X$, el ideal cerrado
$\mathcal K(X)\triangleleft\mathcal B(X)$ conduce al cociente de Calkin. La invertibilidad de
$A+\mathcal K(X)$ equivale a que $A$ sea Fredholm; esta equivalencia es un resultado de teoría de
operadores que se usa después de construir el cociente, no un teorema elemental sobre toda
álgebra de Banach. En el capítulo de investigación se muestran los operadores de Haseman.
